% ---
% id: statistics/binomial-box/def:construction
% kind: def
% title: "The construction map"
% summary: >-
%   The construction map bounds each constrained outcome by its guarded count
%   endpoints divided by the round count, gives every other outcome the bounds
%   0 and 1, and intersects the box with the probability simplex. Its output
%   is a polytope in canonical box form built from the honest envelope.
% prominence: standard
% status: implemented
% uses: [def:endpoint-rule, def:interval, statistics/def:taxonomy, statistics/def:envelope]
% mentions: [req:rule-in-force]
% assumes: []
% refs: []
% provenance: newly-derived
% implements: certiqs_key.statistics.construction:construct
% lean:
% ---

The map is defined for every \tref{admissible endpoint rule}; which rule is
used, by default the \tref{exact binomial quantile rule} of
Definition~\ref{def:interval}, is Requirement~\ref{req:rule-in-force}.

\begin{definition}\label{def:construction}
Let
$(\Chatset,{\honest}^{-},{\honest}^{+},n,\ecomp,\Chatcon,\varepsilon^{-},\varepsilon^{+})$
be the data of an \tref{admissible input} (Definition~\ref{def:endpoint-rule}),
where $({\honest}^{-},{\honest}^{+})$ is its \tref{honest envelope}
(Definition~\ref{def:envelope}) --- an input $\honest$ in the
\tref{point form} being read as the \tref{degenerate envelope}
$(\honest,\honest)$ --- and let $R$ be an \tref{admissible endpoint rule}
(Definition~\ref{def:endpoint-rule}). With the \tref{guarded endpoints}
$k^{\pm}$ of $R$, \eqref{eq:guard}, define
\begin{equation}\label{eq:boxbounds}
  q^{-}(\hat c)\coloneqq
  \begin{cases}k^{-}(\hat c)/n,&\hat c\in\Chatcon,\\[2pt] 0,&\hat c\notin\Chatcon,\end{cases}
  \qquad
  q^{+}(\hat c)\coloneqq
  \begin{cases}k^{+}(\hat c)/n,&\hat c\in\Chatcon,\\[2pt] 1,&\hat c\notin\Chatcon.\end{cases}
\end{equation}
The \term[The map that turns an admissible input and an admissible endpoint
rule into the acceptance box: the guarded count endpoints over n on the
constrained outcomes, 0 and 1 elsewhere, intersected with the
simplex.]{construction map} $\mathsf{Acc}_R$ --- written $\mathsf{Acc}$
where the rule is clear or is the default --- returns the
$\mathrm{polytope}$ geometry in \tref{canonical box form} with these
$(\qlo,\qhi)$, i.e.\ the data \eqref{eq:boxrows}, and
\begin{equation}\label{eq:accset}
  \acSet\ \coloneqq\ \{\mathbf{q}\in\simplex{\Chatset}:\
  q^{-}(\hat c)\le q(\hat c)\le q^{+}(\hat c)\ \ \forall\hat c\in\Chatset\}.
\end{equation}
\end{definition}
